\documentclass[10pt,a4paper]{amsart}
\usepackage[margin=19mm]{geometry}
\usepackage{amsmath,amssymb,mathtools}
\usepackage[scr=rsfs,cal=euler]{mathalfa}
\usepackage{xfrac}
\usepackage[unicode,hidelinks,bookmarks=false]{hyperref}
\newcommand{\ie}{i.e.\ }
\newcommand{\cf}{cf.\ }
\newcommand{\resp}{respectively\ }
\newcommand{\Subsection}{\subsection}
\providecommand{\nobreakdash}{\nobreak\hbox{-}}
\setlength{\emergencystretch}{2em}
\multlinegap0pt
\usepackage{amssymb}
\usepackage{float}
\newtheorem{prop}{Proposition}
\title{Subperiodic groups and bounded automorphisms of periodic graphs}
\author{Igor A. Baburin}
\date{Source: arXiv:2605.12630v1 (CC BY 4.0)}
\begin{document}
\maketitle

\noindent\emph{Source and licence.} Subperiodic groups and bounded automorphisms of periodic graphs. Version arXiv:2605.12630v1 (CC BY 4.0), distributed by arXiv under the Creative Commons Attribution 4.0 International licence, http://creativecommons.org/licenses/by/4.0/, as recorded in the arXiv licence metadata for this exact version and shown on the abstract page https://arxiv.org/abs/2605.12630v1. Full licence notice: \texttt{licenses/arxiv-2605.12630-CC-BY-4.0.txt}.\par\smallskip

\noindent\textit{Editorial scope.} Counted unit: the article's prop prop4 (source0.tex lines 567--570). The article's complete original proof of it is delivered with it (source0.tex lines 572--573, one \texttt{proof} environment of 2 lines). No mathematics is written, repaired or re-derived: the statement and the proof are the article's own lines, in the article's own notation, and no retained line is substituted or re-flowed.  Retained uncounted as the source-local context the proof needs: lines 539--542.  Nothing was omitted from any retained span, so the package carries no omission record.  No retained line contains a citation, so the wrapper carries no bibliography and no bibliography entry.  No retained line contains a figure environment, an image-inclusion command or drawing code, so no image is delivered and the package declares no asset dependency.  Editorial formatting only, in addition: an AMS \texttt{amsart} preamble that restates the article's own declarations and macros used by the retained lines, namely the environments \texttt{prop} on separate counters, together with the standard packages those lines need.  The retained environments print in this document as Proposition 1 in order of appearance, whereas the article prints the same items with the numbering of its own full document.  The complete native source of the article as distributed in the arXiv e-print for this exact version is bundled unchanged as \texttt{sources/200432/source0.tex}.
\begin{prop}
\label{prop2}
Let $\Gamma$ be a vertex-transitive locally finite graph, and let $G < Aut(\Gamma)$ be a regular automorphism group of $\Gamma$. Assume that $C_{Aut(\Gamma)}(G)$ is nontrivial. Put $z \in C_{Aut(\Gamma)}(G)$. Then the generating set $S = \{g | d(x, xg) = 1, x \in V(\Gamma), g \in G \}$ of $G$ is stabilized by some $hz \in G$, where $h \in Stab_{Aut(\Gamma)}(x)$. %, h \neq 1$.
\end{prop}

\begin{prop}
\label{prop4}
Let $\Gamma$ be a Cayley graph of a space group $G$ with respect to the generating set $S = S^{-1}$, and let $Aut(\Gamma)$ be isomorphic to a space group. Then $S$ is stabilized by some $hz \in G$, where $h \in Stab_{Aut(\Gamma)}(x)$ and $z \in C_{Aut(\Gamma)}(G)$. Moreover, $hz$ is either the identity or an element of infinite order.
\end{prop}

\begin{proof} Since $G$ is a subgroup of $Aut(\Gamma)$ of finite index, $z \in C_{Aut(\Gamma)}(G)$ is either the identity or a nontrivial translation. As in Proposition~\ref{prop2}, $S = h^{-1}Sh$ for some $h \in Stab_{Aut(\Gamma)}(x)$. The element $hz$ stabilizes $S$ in the same way $(hz)^{-1}S(hz) = h^{-1}Sh$, if and only if $hz = zh$. If $C_{Aut(\Gamma)}(G) = \{e\}$, then $hz = 1$  since $h = 1$ must hold ($G$ is a regular group of automorphisms). So assume $z \neq 1$. A nontrivial stabilization of $S$ is achieved whenever $h \neq 1$. Then, $hz=zh$ is only possible if $z$ belongs to the invariant subspace of $h$, in which case $hz$ is of infinite order.
\end{proof}

\end{document}
